% Part: first-order-logic
% Chapter: sequent-calculus
% Section: soundness-identity

\documentclass[../../../include/open-logic-section]{subfiles}

\begin{document}

\olfileid{fol}{seq}{sid}

\olsection{السلامة مع \usetoken{S}{identity}}

\begin{prop}
حساب~$\Log{LK}$ المزود بالتتابعيات الابتدائية وقواعد الهوية سليم.
\end{prop}

\begin{proof}
أولًا، التتابعيات الابتدائية من الصورة ${} \Sequent \eq[t][t]$ صالحة:
فلكل !!{structure}~$\Struct M$ يتحقق $\Sat{M}{\eq[t][t]}$.
والحد~$t$ مفروض مغلقًا، أي بلا متغيرات؛ فلا مدخل هنا لإسنادات المتغيرات.

افترض أن الاستدلال الأخير في !!a{derivation} هو~$=$. عندئذ تكون المقدمة
$\eq[t_1][t_2], \Gamma \Sequent \Delta, !A(t_1)$، وتكون النتيجة
$\eq[t_1][t_2], \Gamma \Sequent \Delta, !A(t_2)$. اعتبر
!!a{structure}~$\Struct M$. علينا أن نبين أن النتيجة صالحة، أي إنه إذا
كان $\Sat{M}{\eq[t_1][t_2]}$ و$\Sat{M}{\Gamma}$، فإما أن
$\Sat{M}{!C}$ لبعض $!C \in \Delta$، أو أن $\Sat{M}{!A(t_2)}$.

المقدمة صالحة بفرض الاستقراء. فمتى تحقق
$\Sat{M}{\eq[t_1][t_2]}$ و$\Sat{M}{\Gamma}$، فإما (أ) أن يوجد
$!C \in \Delta$ مع $\Sat{M}{!C}$، وإما (ب) أن يتحقق
$\Sat{M}{!A(t_1)}$. أما (أ) ففيه المطلوب.
وأما (ب)، فخذ~$s$ إسنادًا بحيث $s(x) = \Value{t_1}{M}$.
تعطي \olref[syn][ass]{prop:sentence-sat-true}
$\Sat{M}{!A(t_1)}[s]$. ومن~$\varAssign{s}{s}{x}$، مع
\olref[syn][ext]{prop:ext-formulas}، يلزم
$\Sat{M}{!A(x)}[s]$.
ثم إن~$\Sat{M}{\eq[t_1][t_2]}$ يقتضي
$\Value{t_1}{M} = \Value{t_2}{M}$، فيكون
$s(x) = \Value{t_2}{M}$.
فإذا أعدنا تطبيق \olref[syn][ext]{prop:ext-formulas} حصل
$\Sat{M}{!A(t_2)}[s]$، وبـ
\olref[syn][ass]{prop:sentence-sat-true} نحصل على~$\Sat{M}{!A(t_2)}$، وهو المطلوب.
\end{proof}

\end{document}
